\documentclass[a4paper,11pt]{ltjsarticle}
\usepackage{amsmath,amssymb}
\usepackage{enumitem}
\usepackage{geometry}
\geometry{margin=25mm}

\title{党勢の趨勢のモデル}
\author{}
\date{2026年10月}


\begin{document}
\maketitle

本稿では、党勢の趨勢を推定する現在のモデルを定義する。検討の経緯と、採らなかった案の比較は\texttt{state\_space\_proposal.tex}にある。実装は\texttt{src/seiryoku/state\_space\_glm.py}である。

\section{目的}

モデルの目的は、国政での党勢の動きを、国政選挙と地方選挙の結果から説明することである。測りたいのは国政での党勢であり、地方選挙は数年に一度しかない国政選挙の間を埋める代理の観測と位置づける。国政選挙の結果を事前に当てることは目的としない。ある選挙で期待される結果と実際の結果の差は、その選挙が党勢に与えた衝撃とみなし、党勢の変化をその衝撃の積み重ねとして記述する。

\section{記号}

本節では、以下で使う記号をまとめて定義する。モデルは政党ごとに別々に立てるので、党勢だけでなく水準差や地盤も政党ごとに異なる。添字$p$は政党を表す。ただし党勢$x,z$は、明らかに政党ごとの量なので添字を省く。

\begin{itemize}
  \item $x(t)$: 時刻$t$での政党の党勢。衆院の議席占有率(無所属を除いて正規化)の尺度で測る。
  \item $z(t)=\operatorname{logit}x(t)=\log\{x/(1-x)\}$: logitの空間での党勢。
  \item $\ell$: 選挙の階層(市区町村議会、市区町村長、都道府県議会、都道府県知事、参院選など)。衆院選を基準の階層とし、$\alpha_{p,\text{衆}}=0$とする。
  \item $\alpha_{p\ell}$: 政党$p$の、衆院選に対する階層$\ell$の占有率の水準差。階層$\ell$のすべての自治体に共通である。
  \item $m$: 自治体。$u_{p\ell m}$は、政党$p$の階層$\ell$での自治体$m$に固有の水準(地盤)で、$\alpha_{p\ell}$からのずれを表す。地盤は議会の階層(市区町村議会、都道府県議会)にだけ置き、同じ市でも市議会と県議会の地盤は別の量とする。
  \item 選挙$e$について、$n^\ast_e$は無所属を除いた議席数(首長選では1)、$k_{pe}$はそのうち政党$p$の議席数である。$n^\ast_e$は政党によらない。複数の政党が推薦した首長選(相乗り)では、当選者1人を推薦した政党で分けるので、$k_{pe}$は割合になる。分け方は\ref{sec:coalition}節で述べる。
  \item $v_p$: 政党$p$の直近の国政選挙での比例代表の得票率。
\end{itemize}

\section{状態方程式}

状態は、党勢$z$と、階層ごとの水準差$\alpha_{p\ell}$である。$z$は傾きのないランダムウォーク
\begin{equation}
  dz = \sqrt{q}\,dW
\end{equation}
に従い、$\alpha_{p\ell}$は時間によらない定数とする。分散$q$は全政党に共通の1つで、階層にも時間にもよらない。政党ごとの動きやすさの違いは負荷$\gamma_p$と観測の側($\alpha_{p\ell}$、$\tau^2_{p\ell}$)が受け持つ。$W$は標準ブラウン運動、時間の単位は年である。傾きを持つ形や、水準が跳ぶ形も検討したが、傾きはデータから検出できなかった。跳びは、立憲民主党の2026年衆院選の急落(\ref{sec:shu}節)を含めても対数尤度の上がりが各党1.7以下で、国政選挙の補間も変わらなかった。

\section{観測方程式}

本節では、選挙の種類ごとに、観測が状態とどう結びつくかを定める。国政選挙は$z$をほぼ直接観測し、地方選挙は議席数の分布を通して間接的に観測する。

\subsection{衆院選}
\label{sec:shu}

衆院選は$z$を、議席数の偶然のぶれだけを誤差として観測する。衆院選$e$での政党$p$の議席数$k_{pe}$と、無所属を除いた議席数$n^\ast_e$に対し、
\begin{equation}
  \operatorname{logit}\pi_{pe} = z + \varepsilon,\qquad
  \pi_{pe}=\frac{k_{pe}+0.5}{n^\ast_e+1},\qquad
  \varepsilon\sim N\!\left(0,\ \frac{1}{k_{pe}+0.5}+\frac{1}{n^\ast_e-k_{pe}+0.5}\right)
\end{equation}
とする。誤差の分散は、議席数を$\operatorname{Bin}(n^\ast_e,\ x)$とみなしたときの経験logitの分散で、定数を含まない。0.5議席は、議席0の政党でもlogitを有限にするための補正である。占有率の尺度に直した誤差の標準偏差は、自由民主党の2026年衆院選(315議席)で0.022、参政党(15議席)で0.008、いのちの党(1議席)で0.003である。

誤差を占有率の尺度で全政党に共通の定数$\sigma_{\text{国}}$とする形は採らない。$\sigma_{\text{国}}=0.005$と小さく固定すると、補間の$\pm2$標準偏差に入る割合が0.74にとどまる。尤度から$\sigma_{\text{国}}=0.026$と決めると、議席率が0.03前後の政党では国政選挙がほとんど情報を持たなくなり、補間が地方選挙に引っ張られる(参政党の2026年衆院選で、実際の0.033に対して補間0.195)。二項分布の誤差は、大きな政党では尤度から決めた値に近く、小さな政党では議席数に応じて小さくなる。

2026年衆院選では、立憲民主党と公明党が統一名簿の中道改革連合の名義で戦った。この名義の49議席は、当選者の出身で立憲民主党に21、公明党に28と分けて、それぞれの党の観測とする\footnote{Wikipedia「中道改革連合」(2026年10月18日閲覧):「党内の比率は21議席(立憲系)対28議席(公明系)」。}。地方選挙では合流後も両党が元の名義で戦っているので、合流した名義での推薦は、両党の相乗りとして扱う(\ref{sec:coalition}節)。

\subsection{参院選}

参院選は、衆参それぞれの議席と得票の関係を得票率でつないで観測する。院ごとに、議席占有率$x$と比例代表の得票率$v$の関係
\begin{equation}
  \operatorname{logit}x = k_{\text{院}}\operatorname{logit}v + c_{\text{院}}
\end{equation}
を当てはめる。議席数を$\operatorname{Bin}(n^\ast,\ \operatorname{logit}^{-1}(k_{\text{院}}\operatorname{logit}v + c_{\text{院}}))$とみなし、全政党・全選挙をまとめて$k_{\text{院}},c_{\text{院}}$を最尤推定する。議席数で重みづけるので、議席の少ない党のlogitの大きなぶれに引きずられず、0議席の党も扱える。2017年以降の衆院選で$k_{\text{衆}}=2.13$、2016年以降の参院選の改選議席で$k_{\text{参}}=1.55$である。切片を政党ごとに置くと、選挙を1回ずつ除いて当てる誤差が大きくなる(衆院選で0.029から0.057)。1党あたりの点が2〜4個しかなく、切片が手元の点に合わせすぎるためである。2つの式から$v$を消去すると
\begin{equation}
  \operatorname{logit}x_{\text{参}} = \frac{k_{\text{参}}}{k_{\text{衆}}}\,z + \alpha_{p,\text{参}},\qquad
  \alpha_{p,\text{参}} \sim N\!\left(c_{\text{参}}-\frac{k_{\text{参}}}{k_{\text{衆}}}c_{\text{衆}},\ s^2_{\text{参}}+\left(\frac{k_{\text{参}}}{k_{\text{衆}}}\right)^2 s^2_{\text{衆}}\right)
\end{equation}
となる。$s^2_{\text{院}}$は各院の当てはめの残差分散である。参院選は各党1〜2回しかなく$\alpha_{p,\text{参}}$がデータから決まらないので、その事前分布を2つの曲線から与える。観測誤差は衆院選と同じである。

\subsection{地方選挙}

地方選挙は、無所属を除いた議席のうち政党が得た議席数を、次の2段階で生じるものとして観測する。
\begin{align}
  &\text{地盤:} &&\operatorname{logit}\mu_{pe} = \eta_{pe} = \gamma_p\,z(t_e) + \alpha_{p\ell(e)} + u_{p\ell(e)m(e)},\qquad u_{p\ell m}\sim N(0,\ \tau^2_{p\ell}) \label{eq:stage1}\\
  &\text{議席を数える偶然:} &&k_{pe}\sim\operatorname{Bin}(n^\ast_e,\ \mu_{pe}) \label{eq:local}
\end{align}

1段目は、その自治体で期待される占有率$\mu_{pe}$を決める。$u_{p\ell m}$は自治体$m$の地盤で、選挙をまたいで残る。平均を0とするので、全国に共通のずれは$\alpha_{p\ell}$が持ち、$u_{p\ell m}$は自治体どうしの差だけを表す。$\tau_{p\ell}$はその標準偏差、つまりlogitの尺度での地盤の差の大きさである。首長選では$u_{p\ell m}=0$とする。

2段目は、議席数を$\mu_{pe}$の二項分布で引く。首長選($n^\ast_e=1$)ではベルヌーイ分布になる。$k_{pe}$が割合の首長選では、ベルヌーイ分布の尤度を割合のまま使う。推薦・支持の名簿が公開される前の市区町村長選は、観測に使わない。名簿が出るまでは公認候補が勝った選挙でしか政党の議席が記録されず、結果で選別された観測になるからである。2025年12月22日より後の356件がこれに当たり、名簿が公開されれば通常の市区町村長選として加わる。

占有率ではなく議席数を観測にするのは、次の2つの理由による。

\begin{itemize}
  \item 占有率がちょうど0や1になる議会が大半を占める。議員の大半は無所属なので$n^\ast_e$の中央値は3で、市区町村議会のうち占有率0の議会は自由民主党で70\%、立憲民主党で77\%、国民民主党で89\%、参政党で85\%ある。占有率をベータ分布で観測すると、0や1で尤度が発散するので内側へずらす補正が要るが、そのずらし幅はデータから決まらない。
  \item 同じ占有率でも、議席数によって情報の量が違う。占有率のベータ分布では精度が議席数によらず、2議席の議会の50\%と40議席の議会の50\%を区別できない。二項分布では$n^\ast_e$に応じたぶれが決まり、$n^\ast_e$が小さい議会は自然に情報の少ない観測になる。
\end{itemize}

地盤を入れる理由は、同じ自治体の前回と今回の占有率の相関が0.8前後と高く、自治体間のばらつきの多くが選挙をまたいで残る差だからである。データは2018年から始まるので、2022年以降の地方選挙にはすべて前回の選挙がある。$u_{p\ell m}$を入れることで、実質的に同じ自治体の前回からの変化を観測に使うことになる。地方選挙のばらつきの大半は、地盤と議席を数える偶然で説明される。地盤の差は自由民主党で特に大きく、市区町村議会で$\tau_{p\ell}=1.8$である(日本共産党0.76、立憲民主党0.78)。

採らなかった形は次のとおりである。いずれも国政選挙の補間(\ref{sec:valid}節)が変わらないので、パラメータの少ない形を採った。
\begin{itemize}
  \item 首長選にも地盤を置く形。首長選の$\tau_{p\ell}$は大半の政党で0.05以下と推定され、補間の二乗平均誤差は36件のすべてでほぼ同じだった。
  \item その選挙限りのぶれを精度$\phi_{p\ell}$のベータ二項分布で加える形。対数尤度の上がりは各党1.9以下だった。
\end{itemize}

\subsection{地方選挙の負荷}

負荷$\gamma_p$は、地方の占有率が党勢の変化にどれだけ反応するかを表し、政党ごとにデータから計算する。地方の占有率$\mu$の変化は得票率$v$の変化に等しい($d\mu=dv$)と見立て、衆院の議席と得票の関係と合わせると、
\begin{equation}
  \gamma_p = \frac{d\operatorname{logit}\mu}{dz} = \frac{v_p(1-v_p)}{k_{\text{衆}}\,\bar\mu_p(1-\bar\mu_p)}
\end{equation}
となる。$\bar\mu_p$は推定時点より前の地方議会の議席の合計から求めた占有率である。$v_p$は衆参いずれかの直近の国政選挙の値で、得票率を総務省の集計表から読み取れなかった選挙では、その院の曲線上の点とみなして議席から逆算する。

地方での占有率が得票率より低い政党ほど$\gamma_p$は大きく、地方での変化は国政に小さく写される。2026年9月末時点では、日本共産党0.11、公明党0.19、自由民主党0.49、立憲民主党0.68、日本維新の会1.18、参政党1.47、国民民主党1.50、いのちの党2.02である。

$\gamma_p$を尤度から推定する形も検討したが、国政選挙の補間が悪くなるので採らない。全政党で推定すると補間の二乗平均誤差は0.081から0.085に、尤度の上がりが大きい立憲民主党と日本共産党だけで推定すると0.086になる。立憲民主党の尤度の上がり(12.5)は、$\gamma_p$を0.12まで下げて地方選挙の影響を切り、2026年衆院選の急落(\ref{sec:shu}節)に$z$を合わせた結果で、他の選挙の補間が悪くなる。$d\mu=dv$の代わりに$d\mu/\mu=dv/v$と見立てた$\gamma_p=(1-v_p)/\{k_{\text{衆}}(1-\bar\mu_p)\}$は、全データの対数尤度の合計が8.7下がるので、補間を調べる前に外した。

日本共産党は、この$\gamma_p$を尤度が退ける唯一の政党である。$\gamma_p$を0.35〜0.6にすると対数尤度は4.4上がる。共産党は地方の占有率($\bar\mu_p=0.23$)が得票率($v_p=0.04$)よりずっと高く、地方の変化を得票率の同じ幅の変化とみなす見立てが合わない。地方の変化は$1/\gamma_p\approx9$倍に拡大されて$z$に写るので、共産党の国政選挙の間の動きは大きく読まない。

\subsection{相乗りの首長選}
\label{sec:coalition}

相乗りの首長選では、当選者1人を、選挙直前の各党の強さに比例して推薦した政党に分ける。元のデータは推薦した政党で等分しているが、等分では強い党と弱い党の相乗りでも同じだけ数えることになる。相乗りの首長選は397件あり、そのうち243件は自由民主党と公明党の組である。

分け方は次の2段階で決める。
\begin{enumerate}
  \item 等分のデータで、各党のモデルのパラメータを推定する。
  \item 相乗りの首長選$e$ごとに、推薦した政党の組$S_e$について、各党のフィルタが選挙直前に予測する当選確率$\hat\mu_{pe}=\operatorname{logit}^{-1}(\hat\eta_{pe})$を求め、組の取り分の合計を
  \begin{equation}
    k_{pe} = \Bigl(\sum_{q\in S_e}k^{\text{等分}}_{qe}\Bigr)\frac{\hat\mu_{pe}}{\sum_{q\in S_e}\hat\mu_{qe}}
  \end{equation}
  で分け直す。$S_e$は推定の対象とする政党に限り、それ以外の政党の取り分は等分のままとする。
  \item 分け直したデータで、パラメータを推定し直す。
\end{enumerate}
増える定数はない。分け直しを3回繰り返しても1回の場合とほとんど変わらないので、1回にとどめる。

\section{パラメータの決め方}

モデルの定数は、データから推定するもの、推定の外のデータから計算するもの、結果に効かないことを確かめたものに分かれる。値をデータから決められない定数に依存する形は採らない。

\begin{table}[h]
\centering
\begin{tabular}{lll}
\hline
定数 & 決め方 & 数 \\
\hline
$q$ & 全政党の対数尤度の合計を最大化 & 全体で1 \\
$\tau^2_{p\ell}$ & 政党・議会の階層ごとに最尤推定 & 2 \\
$k_{\text{衆}}, c_{\text{衆}}, k_{\text{参}}, c_{\text{参}}, s^2_{\text{衆}}, s^2_{\text{参}}$ & 総務省の公表値から計算 & --- \\
$\gamma_p$ & 議席と得票の関係と地方の議席から計算 & --- \\
状態の事前分散 & 10。100にしても結果がほぼ変わらない & --- \\
\hline
\end{tabular}
\end{table}

$q$は、いくつかの値で各党の$\tau^2_{p\ell}$を最尤推定して対数尤度の合計を求め、最大点の付近を$\log q$の2次式で補間して決める。2026年9月末時点で$q=0.024$である。$q$を政党ごとに推定する形と比べて、国政選挙の補間の二乗平均誤差は変わらず(どちらも0.081)、パラメータは7つ減る。政党ごとの形では、日本共産党の分散が0.49と他党(0.003〜0.027)よりはるかに大きく推定され、2020年前後の$\pm2$標準偏差の帯が0.02から0.59まで広がっていた。

\section{推定の計算}

本節では、与えたパラメータのもとで状態を推定する手順を述べる。観測が状態の非線形関数なので、地方選挙ごとに正規分布で近似してカルマンフィルタで更新する。

\begin{enumerate}
  \item 地方選挙$e$の直前の状態の分布から、$\eta_{pe}$の予測分布$N(\hat\eta_{pe}, s^2_{pe})$を求める。$u_{p\ell m}$の分布は、同じ階層のその自治体の前回の選挙での事後分布を使う(初めての自治体では$N(0,\tau^2_{p\ell})$)。
  \item 式\eqref{eq:local}の尤度と予測分布の積を$\eta_{pe}$について最大化し、その点での2階微分から事後分布を正規分布で近似する(ラプラス近似)。
  \item 同じ事後分布を与える正規の擬似観測$\tilde y_{pe}$とその分散$R_{pe}$を求め、カルマンフィルタで状態と$u_{p\ell m}$を更新する。$u_{p\ell m}$は状態ベクトルには入れず、階層と自治体の組ごとに、事後分布の平均と分散、状態との共分散を覚えておく。捨てるのは自治体どうしの共分散だけである。模擬データでは、$u_{p\ell m}$を状態に入れた厳密な計算と比べて、$z$の推定値の差は事後標準偏差の2\%にとどまる。状態との共分散まで捨てると、この差は15\%になり、$z$の推定誤差も1割大きくなる。
  \item 対数尤度も同じ近似で求める。事前分布の影響を避けるため、最初の2つの選挙と、各階層で最初の選挙は尤度に含めない。
\end{enumerate}

過去の時点の状態は、その後の選挙も使って推定する。前向きのフィルタの結果から、固定区間の平滑化(RTS)で各時点の状態の平均と分散を求める。地方選挙はフィルタでの擬似観測のまま扱う。

\section{出力}

月ごとに、次の2つを出力する。

\begin{enumerate}
  \item 月末の党勢$\hat x=\operatorname{logit}^{-1}\hat z$と、その幅$\operatorname{logit}^{-1}(\hat z\pm2\hat s)$。$\hat z,\hat s$は平滑化の平均と標準偏差である。平滑化は後の選挙も使うので、大きな国政選挙での跳びを前の月に割り振り、選挙の前から党勢が動いていたように見せる。そこで、その月までの選挙だけによるフィルタの平均と標準偏差による値も併せて出力する。最新の月では両者は一致する。
  \item その月の選挙ごとの寄与。
\end{enumerate}

選挙ごとの寄与は、その月の党勢の推定値の変化を選挙ごとにちょうど分けたものである。傾きを置かないので、前向きのフィルタでは予測の段階で$z$の平均は動かず、$z$は観測ごとの更新でだけ動く。選挙$e$による$z$の変化を$\Delta z_e$とすると、ある月の$\hat z$の変化はその月の$\Delta z_e$の和に等しい。各選挙について、期待される占有率$\operatorname{logit}^{-1}(\hat\eta_{pe})$、実際の占有率$k_{pe}/n^\ast_e$、$\Delta z_e$を出す。$\Delta z_e$は、期待と実際の差、つまりその選挙の衝撃に、その選挙の情報の重みを掛けたものである。

\section{検証}
\label{sec:valid}

モデルが党勢の動きを説明できているかは、国政選挙の補間で確かめる。国政選挙を1回ずつ除き、残りの全期間のデータから平滑化でその結果を再現できるかを見る。除いた選挙の議席と得票の点は、$\gamma_p$、参院選のつなぎ、$q$の計算からも除く。比較の基準は、同じモデルから地方選挙を除いた、国政選挙だけの補間である。

モデルの形も、この補間で選んだ。2つの形の補間の二乗誤差の差を件ごとに求め、その平均が標準誤差の2倍に届かなければ、パラメータの少ない形を採る。同じ選挙の件どうしは相関するので、標準誤差は選挙ごとに差を合計してから求める。全データの対数尤度は、補間を調べる前に明らかに悪い形を外すためにだけ使った。

各党の最初の国政選挙を除く36件の結果を表に示す。

\begin{table}[h]
\centering
\begin{tabular}{lcccc}
\hline
 & 二乗平均誤差 & 誤差の中央値 & $\pm2$標準偏差に入る割合 & 標準化誤差の標準偏差 \\
\hline
国政選挙だけ & 0.098 & 0.023 & 0.89 & 1.6 \\
本稿のモデル & 0.081 & 0.021 & 0.83 & 1.7 \\
\hline
\end{tabular}
\caption{国政選挙の補間の誤差(36件)。}
\end{table}

地方選挙を使うと、二乗平均誤差は0.098から0.081に、誤差の中央値は0.023から0.021に下がる。地方選挙は、国政選挙の間の党勢の動きの一部を説明している。政党別の絶対誤差の平均は表のとおりで、改善は参政党、自由民主党、立憲民主党で大きい。日本維新の会、いのちの党、公明党では、国政選挙だけの補間がすでに0.01〜0.02と小さく、地方選挙を加えるとわずかに悪くなる。

\begin{table}[h]
\centering
\begin{tabular}{lccc}
\hline
政党 & 件数 & 国政選挙だけ & 本稿のモデル \\
\hline
自由民主党 & 5 & 0.174 & 0.131 \\
立憲民主党 & 5 & 0.152 & 0.109 \\
参政党 & 3 & 0.059 & 0.035 \\
国民民主党 & 5 & 0.030 & 0.025 \\
公明党 & 5 & 0.016 & 0.017 \\
日本維新の会 & 5 & 0.015 & 0.025 \\
いのちの党 & 3 & 0.013 & 0.019 \\
日本共産党 & 5 & 0.009 & 0.008 \\
\hline
\end{tabular}
\caption{政党別の補間の絶対誤差の平均。}
\end{table}

大きな外れは、地方選挙にも前後の国政選挙にも現れなかった国政の動きである。自由民主党の2026年衆院選(実際0.686、補間0.395)と2021年衆院選(0.577、0.728)、立憲民主党の2026年衆院選(0.046、0.263)と2024年衆院選(0.330、0.146)、参政党の2025年参院選(0.124、0.036)がこれに当たる。立憲民主党の2024年衆院選は、国政選挙だけの補間でも0.122で、2022年と2025年の参院選から見込まれる水準を大きく上回った。

この補間の成績を見ながらモデルの形を選んでいるので、成績はいくらか楽観的に出ている。選択の偏りを抑えるため、差が標準誤差の2倍に届かなければパラメータの少ない形を採り、比べる形の数もこれ以上増やさない。

相乗りの首長選の分け直し(\ref{sec:coalition}節)は、補間の成績をほとんど変えない。中道改革連合を分ける前の版では、等分のままで二乗平均誤差0.067、誤差の中央値0.019、$\pm2$標準偏差に入る割合0.82で、分け直すと0.068、0.019、0.76になる。入る割合の差は、区間の端から0.003以内にあった2件(自由民主党の2021年衆院選、国民民主党の2026年衆院選)が外れたことによる。分け直しを採るのは補間の成績のためではなく、等分では強い党と弱い党を同じだけ数えてしまうからである。

\section{国政だけの成分}

地方選挙に現れない国政だけの党勢の動きは、本体のモデルには入れず、診断用の出力とする。国政の党勢を$z+z_N$とし、地方選挙の負荷は$z$にだけかける。$z_N$は始点で0、分散$q^{\mathrm{N}}_p$のランダムウォークで、$q^{\mathrm{N}}_p$は尤度から推定する。

本体に入れない理由は、国政選挙の補間が悪くなるからである。相乗りを等分していた時点の検証(34件)では、$z_N$を加えると全期間の対数尤度は自由民主党で7.1、国民民主党で9.5、参政党で7.7、立憲民主党で2.9上がった(パラメータは各党1つ増える)。一方、補間の二乗平均誤差は0.069から0.080に悪化した。$z_N$は国政選挙でしか動かないので、除いた選挙の値を地方選挙ではなく前後の国政選挙から決めるようになるためである。悪化の大半は、$q^{\mathrm{N}}_p$が大きく推定される参政党で生じた。

診断用には、$z_N$の平滑化の平均と標準偏差を月ごとに出力する。$z_N$が大きく0から離れる時期は、国政選挙の結果が地方選挙から期待される水準から外れた時期である。

\end{document}
